\paragraph{A rule that lets incumbents move}
For each training lottery compute the blind allocation. In grades above pre-kindergarten, keep every old proposal up through the contract the pupil holds (all proposals for an unplaced pupil), and forbid adding a course that would be proposed to before that point. Intersect these tests across draws. The grade-by-grade proposal sequence, including family promotions and current-school fallbacks, is then identical in every higher grade.

For the last grade adopt a weaker test. Require the original course of every pupil placed on any training draw. A newly added course must either (i) fail to block the baseline in every training draw, or (ii) have at least as many ordinary and reserved seats in total as there are potential applicants in the union of the two policies. Hold common courses and individual quotas fixed. The added freedom in (ii) can change the allocated school.

\emph{Claim.} Every baseline recipient still has a seat on every training draw.
Fix one draw. The upper grades finish exactly as before, so the final grade faces unchanged contract priorities, preference orders and family states. For an incumbent assigned to a universally slack course (ii), truncate their new list after \emph{all} types at that course. For every other incumbent truncate immediately after the particular old contract, including its type. Previously unassigned pupils keep their whole new menus. At a non-(ii) pool any applicant proposing on the truncated lists either held a seat in that pool at baseline or ranks the contract above their old one. In the latter case stability on an old course, or the nonblocking test on a newly displayed course, means the pool was full with better-priority holders. This includes applicants who prefer a different type at their original course. The extra seat types allowed at the slack courses introduce no extra proposals at these ordinary pools. An old incumbent who reaches their floor pool is therefore among its highest-priority applicants and secures a seat. The guaranteed current-school fallback cannot reject its holder either. A pupil who reaches a universally slack course cannot be rejected by every type, since that would require capacity-many distinct \emph{other} applicants holding its seats. Thus at least one type on their truncated list accepts them. Every incumbent is therefore assigned on these truncated lists. Restoring options below each truncation cannot change the proposal path: no incumbent exhausted their list of courses in this run. The resulting full lists preserve all incumbent placements. A final-grade incumbent may switch schools, but their new course ranks at least as high in the institution's family-adjusted ordering. A pupil who stays at a slack course may receive a different seat type, and an initially unplaced pupil may receive a seat.
